% d_i is route distance in kilometres and g_i is route gradient in percent.
\begin{align}
\eta_i
 &= \beta_0 + \beta_1 d_i + \beta_2 \sqrt{d_i}
  + \beta_3 d_i^2 + \beta_4 (g_i-g_0) \notag\\
 &\quad + \beta_5 d_i(g_i-g_0)
  + \beta_6 \sqrt{d_i}(g_i-g_0),
\label{eq:pct-predictor}\\
p_i^{\mathrm{PCT}}
 &= \expit(\eta_i)
  = \frac{1}{1+\exp(-\eta_i)}.
\label{eq:pct-probability}
\end{align}

% The coefficient set identifies a named PCT scenario. In this workbench the
% probability is a transferable spatial propensity prior, not an Auckland-
% estimated causal response or forecast.
